\documentclass[10pt]{article}
\usepackage[margin=0.74in]{geometry}
\usepackage{amsmath,amssymb,booktabs,graphicx,microtype}
\graphicspath{{../figures/pdf/}}
\title{Mean-Field Capacity Control for Production-Derived Serverless Workloads}
\author{Koki Matsuoka\\Nagoya University\\Independent Computational Study}
\date{2026}
\begin{document}
\maketitle
\begin{abstract}
We study capacity control for trace-derived serverless demand shapes using finite-$N$ JSQ($d$) simulation and a quasi-stationary mean-field model-predictive controller. Public Azure Functions traces supply workload shapes; all performance figures are controlled simulations, not deployment measurements. Gaussian-process forecasts provide both a point forecast and an uncertainty-aware upper scenario. In the high-volume scenario, EWMA MF-MPC reduced p95 simulated waiting by 56.7\% relative to reactive control (95\% paired bootstrap CI [52.8\%, 60.6\%]; mean resource fractions 0.274 vs 0.273). The paired interval excludes zero. GP-UCB reaches a much lower-delay operating point, but it uses substantially more resources and is not resource-matched; it therefore does not establish equal-cost superiority.
\end{abstract}
\section{Introduction}
Serverless demand can be nonstationary while capacity changes are delayed. This study asks what a mean-field queueing proxy contributes when decisions are constrained by a validation-defined resource budget. The conclusion is deliberately conservative: the strongest strict result is for a simple EWMA forecast, while uncertainty-aware GP-UCB remains a descriptive tradeoff rather than a matched-budget win.
\section{Data and experimental protocol}
The Azure Functions Invocation Trace 2021 provides production-derived workload \emph{shapes} \cite{azuretrace2021}. Request time is reconstructed as end timestamp minus duration, and counts are binned in five-minute intervals. For the revised analysis, data-quality eligibility criteria were fixed before rerunning controller evaluation after the earlier degenerate held-out scenario was discovered. Eligibility is evaluated before testing; among eligible applications, ranking uses training statistics only. GP fitting uses training data, beta calibration and alpha matching use validation only, and test data are used only after settings are frozen.

The simulator rescales intensity to a controlled stress regime. Thus the reported p95 values are simulated queue waiting times, not Azure request latency. Severe overload can leave queues draining for a long time, explaining the large bursty values.
\section{Mathematical model}
\subsection{Finite-$N$ JSQ($d$) system}
There are $N$ parallel queues. Trace-derived arrivals form a time-varying intensity $\lambda(t)$; each arrival samples $d$ queues and joins a shortest sampled queue (JSQ($d$)). Service has mean $1/\mu$; the event-driven finite-$N$ simulator is the experimental ground truth. Let $Q_i(t)$ denote queue $i$ and let $q_k(t)$ denote the fraction with $Q_i(t)\ge k$.
\subsection{Mean-field approximation}
For a stationary rate and exponential service, the supermarket-model motivation is
\begin{equation}
\frac{d q_k}{dt}=\lambda\left(q_{k-1}^d-q_k^d\right)-\mu\left(q_k-q_{k+1}\right),\qquad k\ge1,
\end{equation}
with $q_0=1$. With $\rho=\lambda/(N\mu)$, the stationary tail is $q_k^*=\rho^{(d^k-1)/(d-1)}$ for $d\ge2$, and $q_k^*=\rho^k$ for $d=1$ \cite{mitzenmacher2001,ying2016}. The controller uses this quasi-stationary fixed point, not a transient controlled mean-field ODE.
\subsection{Receding-horizon capacity search}
At each decision time, a discrete candidate capacity $c$ is searched over a horizon of forecast rates. The implemented objective is the sum of predicted quasi-stationary waiting cost normalized by mean service, an $\alpha c/N$ resource penalty, and a switching penalty. Capacity actions are ramp-limited, delayed, and quantized; only the first action is applied before replanning.
\section{Forecasting and uncertainty}
GPs are fitted only on training rows \cite{rasmussen2006}. GP-mean is a point forecast; GP-UCB is an upper demand scenario and is evaluated by coverage, underprediction, and width rather than point-forecast MAE/RMSE.
\begin{table}[!htbp]\centering\small\caption{Point forecast metrics.}\begin{tabular}{lrr}\toprule Method & MAE & RMSE\\\midrule
ewma & 25.3 & 43.8 \\
gp\_mean & 33.8 & 58.2 \\
last\_value & 23.7 & 45.9 \\
seasonal\_naive & 52.5 & 77 \\\\[2pt]\bottomrule\end{tabular}\end{table}
\begin{table}[!htbp]\centering\small\caption{GP-UCB uncertainty metrics.}\begin{tabular}{lrrrr}\toprule Method & Coverage & Underprediction & Log width & Normalized width\\\midrule
GP-UCB & 0.917 & 0.083 & 1.07 & 95.9 \\\\[2pt]\bottomrule\end{tabular}\end{table}
\section{Capacity controllers and resource matching}
Reactive validation resource usage defines the reference budget. For each predictive controller, alpha is selected by adaptive log-space search using validation only. The configured relative tolerance is 2\%. Because capacity is discrete, a match can be impossible; unmatched points stay in Pareto plots but are excluded from strict matched-budget claims.
\section{Experimental results}
\subsection{Scenario quality and primary outcomes}
All three reported scenarios meet the revised split-level activity rules. Figure~\ref{fig:primary} uses separate log-scale panels so that the bursty stress regime does not obscure high-volume or diurnal results.
\begin{figure}[!htbp]\centering\includegraphics[width=\textwidth]{fig08_primary_controller_results.pdf}\caption{Primary paired-seed p95 simulated waiting results (log scale). Each panel is a scenario; values are simulated waiting times, not Azure deployment latency.}\label{fig:primary}\end{figure}
\begin{table}[!htbp]\centering\scriptsize\caption{Primary p95 waiting seconds and validation matching status.}\begin{tabular}{llrrl}\toprule Scenario & Controller & Mean & 95\% CI & Match\\\midrule
bursty & ewma\_mf\_mpc & 4.11e+05 & [3.94e+05, 4.26e+05] & unmatched \\
bursty & gp\_mean\_mf\_mpc & 5.09e+04 & [4.72e+04, 5.51e+04] & unmatched \\
bursty & gp\_ucb\_mf\_mpc & 5.14e+04 & [4.75e+04, 5.59e+04] & unmatched \\
bursty & oracle\_mf\_mpc & 6.5e+04 & [6e+04, 6.97e+04] & unmatched \\
bursty & reactive & 5.41e+05 & [5.22e+05, 5.59e+05] & reference \\
bursty & seasonal\_mf\_mpc & 7.21e+05 & [6.99e+05, 7.39e+05] & unmatched \\
bursty & static & 7.23e+05 & [7.01e+05, 7.43e+05] & not\_matched\_baseline \\
diurnal & ewma\_mf\_mpc & 18.7 & [18.4, 18.9] & matched \\
diurnal & gp\_mean\_mf\_mpc & 22.9 & [22.7, 23.1] & matched \\
diurnal & gp\_ucb\_mf\_mpc & 22.9 & [22.7, 23.1] & matched \\
diurnal & oracle\_mf\_mpc & 22.9 & [22.7, 23.1] & matched \\
diurnal & reactive & 20.3 & [20.2, 20.5] & reference \\
diurnal & seasonal\_mf\_mpc & 22.9 & [22.7, 23.1] & matched \\
diurnal & static & 10.8 & [10.8, 10.9] & not\_matched\_baseline \\
high\_volume & ewma\_mf\_mpc & 1.65e+04 & [1.42e+04, 1.89e+04] & matched \\
high\_volume & gp\_mean\_mf\_mpc & 3.12e+04 & [2.99e+04, 3.26e+04] & matched \\
high\_volume & gp\_ucb\_mf\_mpc & 2.04e+03 & [1.9e+03, 2.18e+03] & unmatched \\
high\_volume & oracle\_mf\_mpc & 2.19e+03 & [1.9e+03, 2.54e+03] & unmatched \\
high\_volume & reactive & 3.87e+04 & [3.27e+04, 4.55e+04] & reference \\
high\_volume & seasonal\_mf\_mpc & 3.07e+05 & [2.99e+05, 3.15e+05] & unmatched \\
high\_volume & static & 602 & [560, 641] & not\_matched\_baseline \\\\[2pt]\bottomrule\end{tabular}\end{table}
\subsection{Strict matched-budget comparison}
The predeclared primary comparison is high-volume EWMA MF-MPC versus reactive control. EWMA MF-MPC reduced p95 simulated waiting by 56.7\% relative to reactive control (95\% paired bootstrap CI [52.8\%, 60.6\%]; mean resource fractions 0.274 vs 0.273). The paired interval excludes zero. The resource mismatch is below the configured 2\% threshold, so this is the principal strict matched-budget result.
Matching status was determined on validation only; reported resource fractions below are test-set means.
\begin{table}[!htbp]\centering\caption{High-volume paired EWMA MF-MPC versus reactive control.}\begin{tabular}{rrrrr}\toprule Reduction & 95\% CI & EWMA resource & Reactive resource & Mismatch\\\midrule
56.7\% & [52.8, 60.6]\% & 0.274 & 0.273 & 0.06\% \\\\[2pt]\bottomrule\end{tabular}\end{table}
\subsection{Pareto tradeoffs and uncertainty-aware control}
GP-UCB MF-MPC reduced p95 simulated waiting by 94.4\% relative to reactive control (95\% paired bootstrap CI [93.5\%, 95.3\%]; mean resource fractions 0.758 vs 0.273). The paired interval excludes zero. At its validation-selected operating point, GP-UCB is therefore descriptive evidence of a lower-delay, higher-resource region, not a matched-resource or equal-cost conclusion.
\begin{figure}[!htbp]\centering\includegraphics[width=\textwidth]{fig09_pareto.pdf}\caption{Validation and test Pareto operating points. Markers distinguish the validation-selected setting from held-out test outcomes; test data do not select alpha.}\end{figure}
\subsection{Robustness}
Bursty delay, population, and service-model checks are retained as stress tests. A compact high-volume activation-delay check evaluates only the predeclared EWMA-versus-reactive comparison. EWMA has lower mean p95 waiting at both checked delays; the uncertainty intervals are shown without converting this robustness check into a new tuned selection.
\begin{table}[!htbp]\centering\scriptsize\caption{High-volume delay robustness; p95 simulated waiting seconds.}\begin{tabular}{rrrrr}\toprule Delay bins & EWMA mean & EWMA 95\% CI & Reactive mean & Reactive 95\% CI\\\midrule
0 & 2.02e+03 & [1.88e+03, 2.16e+03] & 4.93e+03 & [3.88e+03, 6.2e+03] \\
2 & 5.55e+04 & [4.6e+04, 6.55e+04] & 6.52e+04 & [5.38e+04, 7.71e+04] \\\\[2pt]\bottomrule\end{tabular}\end{table}
\begin{figure}[!htbp]\centering\includegraphics[width=0.74\textwidth]{fig13_high_volume_delay.pdf}\caption{High-volume activation-delay robustness for EWMA MF-MPC and reactive control.}\end{figure}
\section{Limitations}
Demand time is reconstructed as end timestamp minus duration. Workload shapes are production-derived but intensities are scaled; the finite-$N$ simulator is not Azure deployment. The quasi-stationary predictor assumes stationary/exponential-service structure, while empirical-service simulation deliberately creates model mismatch. GP fitting can subsample for tractability; only three representative scenarios are studied. Discrete capacity can prevent matching, and unmatched GP-UCB comparisons cannot establish equal-cost superiority.
\clearpage
\section{Conclusion}
The study does not establish a strict equal-resource advantage for uncertainty-aware GP-UCB control. It does show that GP-UCB can move a high-volume simulation to a lower-delay, higher-resource operating region. Under the validation-defined matched budget, the strongest supported finding is the paired high-volume EWMA MF-MPC comparison with reactive control. These claims are simulation-only and trace to machine-readable artifacts.
\section{Bilingual reporting}
The Japanese report and Japanese claim translations are provided in \texttt{paper/report\_ja.md} and \texttt{docs/CLAIMS.md}; numerical values are generated from the same machine-readable artifacts.
\bibliographystyle{plain}
\bibliography{references}
\end{document}
